\documentclass[11pt]{article}
\usepackage[margin=1in]{geometry}
\usepackage{amsmath,amssymb}
\usepackage{booktabs}
\usepackage{listings}
\usepackage{hyperref}
\usepackage{parskip}
\lstset{basicstyle=\ttfamily\small, breaklines=true, columns=fullflexible, frame=single}

\title{Independent Replication Report\\
\large arXiv:quant-ph/9804044 --- \emph{``Quantum Computing''} (Scarani, 1998)}
\author{Ollie / OpenClaw replication wave QC-200}
\date{2026-07-05}

\begin{document}
\maketitle

\section{Paper summary}

\textbf{Paper.} V.~Scarani, \emph{Quantum Computing}, Am.~J.~Phys.~\textbf{66} (11), pp.~956--960 (November~1998); arXiv:quant-ph/9804044v2 (6~Oct~1998), 10 pages.

\textbf{Type.} Single-author pedagogical review / survey aimed at undergraduate physics students.  Not an original algorithm paper --- the stated purpose is to give ``a self-contained description of a quantum computer based on well-known elements of undergraduate quantum mechanics.''

\textbf{Framework Scarani exposits.}  Scarani builds QC out of two ingredients, following the universality theorem of Barenco~\emph{et~al.} (\emph{Phys.~Rev.~A} \textbf{52}, 3457, 1995): (i)~arbitrary one-qubit rotations $R_x, R_y, R_z$ implemented as NMR pulses in a rotating frame, and (ii)~controlled-NOT (CNOT / XOR) gates implemented by uniquely addressing a two-spin transition frequency in the Hamiltonian
\[
H_0 = -\tfrac{\hbar}{2}\!\left[ \Omega_1 (\sigma_z\!\otimes\!\mathbf{1}) + \Omega_2 (\mathbf{1}\!\otimes\!\sigma_z) + \omega_c (\sigma_z\!\otimes\!\sigma_z) \right].
\]

\textbf{Exercises posed.}  Section~3 poses four worked exercises: (3.1)~preparing a GHZ state $\tfrac{1}{\sqrt 2}(|{+}{+}{+}\rangle+|{-}{-}{-}\rangle)$ from $|{+}{+}{+}\rangle$ using $R_y(\pi/4)$ and two CNOTs; (3.2)~the two-qubit NOT gate as a product of two one-spin $R_x(\pi/2)$'s and hence unable to modify entanglement; (3.3)~a Bell-basis readout gate $T = (\mathbf{1}\otimes R_y(\pi/4))\,C^1_{2-}$; (3.4)~the two-qubit Quantum Fourier Transform matrix (eq.~22).

\textbf{Algorithm references.}  Scarani's post-publication (v2) reference list explicitly cites the NMR realization of the Deutsch--Josza algorithm (Chuang, Vandersypen, Zhou, Leung, Lloyd, \emph{Nature}~\textbf{393} (1998) 143), NMR realization of Grover (Jones, Mosca, Hansen, quant-ph/9805069), and Shor's factoring (his ref.~[2]).  His Section 3.4 QFT is the linchpin of Shor's algorithm.

\section{What we replicated}

Because the paper is a survey with no headline numerical claim to reproduce, the QC wave brief specifies exercising two \emph{representative} algorithms from the survey's own frame of reference in Qiskit statevector simulation to verify that Scarani's ``$R_u(\theta)$~+~CNOT'' toolkit does deliver the promised quantum advantage.  We additionally reproduce Scarani's own Exercise~3.4 as a direct numerical cross-check against a matrix he actually writes down.

\begin{enumerate}
    \item \textbf{Deutsch--Jozsa}, $n=3$ (i.e.\ $N=2^n=8$).  A single oracle query must distinguish constant from balanced with success probability~1: if $P(\text{input register} = |0\rangle^{\otimes 3}) = 1$ the function is constant; if $= 0$ the function is balanced.  Referenced by Scarani (NMR realization by Chuang \emph{et~al.}, \emph{Nature} \textbf{393}, 143, 1998).
    \item \textbf{Simon's algorithm}, $n=3$.  For a function $f:\{0,1\}^3\to\{0,1\}^3$ satisfying $f(x)=f(x\oplus s)$ for hidden $s\in\{0,1\}^3\setminus\{0\}$, recover $s$ in $O(n)$ quantum queries plus $O(n^3)$ classical GF(2) linear algebra.  Simon's algorithm is the direct progenitor of the period-finding subroutine at the heart of Shor (Scarani's ref.~[2]).
    \item \textbf{Scarani exercise 3.4 (QFT$_{n=2}$, eq.~22)}, as a bonus reproduction check against a matrix the paper actually prints.
\end{enumerate}

\section{Claims table}

\begin{tabular}{lp{6.5cm}ll}
\toprule
ID & Claim & Testable? & Tested?\\
\midrule
C1 & Universality of one-qubit rotations + CNOT (Scarani citing Barenco \emph{et~al.}~[8]) & Yes (via exercising complete algorithms) & Yes, indirectly (DJ, Simon, QFT all built from $H,X,\text{CNOT}$) \\
C2 & Exercise 3.4: QFT$_{n=2}$ matrix (eq.~22) & Yes (matrix identity) & \textbf{Yes, direct} \\
C3 & Framework supports quantum advantage (implicit; motivates the paper) & Yes (via DJ 1-query and Simon poly-query) & \textbf{Yes, direct} \\
C4 & NMR selectivity conditions $\omega_1{-}\omega_2 > 2\omega_c$ etc. & Physical (not simulable in statevector) & No (out of scope) \\
C5 & Exercise 3.1 GHZ sequence produces $\tfrac{1}{\sqrt 2}(|{+}{+}{+}\rangle+|{-}{-}{-}\rangle)$ & Yes (statevector) & Not run (trivial; DJ/Simon/QFT already cover $H+\text{CNOT}$ toolkit) \\
C6 & Exercise 3.2 NOT = product of two $R_x(\pi/2)$'s cannot modify entanglement & Yes (statevector) & Not run (trivially checked from tensor structure) \\
\bottomrule
\end{tabular}

\section{Method (numbered, exact commands)}

\subsection{Environment}
\begin{lstlisting}
$ python3 -m venv .venv
$ .venv/bin/pip install -q --upgrade pip
$ .venv/bin/pip install -q qiskit qiskit-aer numpy
$ .venv/bin/python -c "import qiskit, qiskit_aer; \
      print(qiskit.__version__, qiskit_aer.__version__)"
2.5.0 0.17.2
\end{lstlisting}

Host: macOS Darwin~25.3.0, Python~3, no GPU used.  Runtime: 2.45~s wall.

\subsection{Deutsch--Jozsa}

Oracle construction (implements $U_f|x\rangle|y\rangle = |x\rangle|y\oplus f(x)\rangle$):
\begin{itemize}
    \item \textbf{Constant $f\equiv 0$:} empty circuit.
    \item \textbf{Constant $f\equiv 1$:} X on the ancilla.
    \item \textbf{Balanced:} draw a nonzero mask $a\in\{0,1\}^n$; for each $i$ with $a_i=1$ place CNOT($i \to$ ancilla).  Then $f(x)=a\cdot x \bmod 2$ is balanced.
\end{itemize}
Circuit: put ancilla into $|-\rangle$, put input register into $|+\rangle^{\otimes n}$, apply $U_f$, Hadamard the input register, read out.  We compute $P(\text{input}=|0\rangle^{\otimes n})$ from the exact statevector (no sampling).

\subsection{Simon's algorithm}

Standard 2$n$-qubit oracle: copy $x$ into $y$ via CNOTs $i\to n+i$, then if $s\ne 0$ let $j$ be the first bit with $s_j=1$ and add controlled-$s$ from qubit $j$ into the second register.  This implements a 2-to-1 $f$ with $f(x)=f(x\oplus s)$.

Per Simon round: apply $H^{\otimes n}$ to the input register, apply oracle, apply $H^{\otimes n}$, measure input.  Any measured $y$ satisfies $s\cdot y \equiv 0 \bmod 2$.  Collect $n{-}1$ linearly independent nonzero equations; solve nullspace over $\mathrm{GF}(2)$; the unique nonzero nullspace vector is $s$.  We ran 5 different hidden $s$ values and always recovered $s$ in $\le 4$ rounds.

\subsection{QFT$_{n=2}$ cross-check}
Construct $F_{kx} = e^{2\pi i k x / 4} / 2$ for $k,x\in\{0,1,2,3\}$ and compare element-wise to Scarani's eq.~(22).

\subsection{Reproducible driver}
All three are wrapped in \verb|report/evidence/run_algorithms.py| with fixed seeds.
Re-run:
\begin{lstlisting}
$ .venv/bin/python report/evidence/run_algorithms.py
{
  "dj_constant_P0": 1.0,
  "dj_balanced_P0": 7.925760013168616e-64,
  "simon_success_rate": 1.0,
  "scarani_qft2_matches": true,
  "verdict_dj": "PASS",
  "verdict_simon": "PASS"
}
\end{lstlisting}

\section{Results vs.\ paper}

\begin{tabular}{p{5.2cm}lll}
\toprule
Test & Expected (paper / theory) & Observed (this run) & Match?\\
\midrule
DJ~$n{=}3$, constant $f{=}0$ & $P(|0^3\rangle_{\text{in}})=1$ & $1.0$ (exact) & \textbf{Yes} \\
DJ~$n{=}3$, constant $f{=}1$ & $P(|0^3\rangle_{\text{in}})=1$ & $1.0$ (exact) & \textbf{Yes} \\
DJ~$n{=}3$, balanced (5 masks) & $P(|0^3\rangle_{\text{in}})=0$ & $\le 8\!\times\!10^{-64}$ (numerical zero) & \textbf{Yes} \\
Simon $n{=}3$, $s\in\{101,011,110,111,100\}$ & Recover $s$ in $O(n)$ quantum rounds & 5/5 recovered in 2--4 rounds & \textbf{Yes} \\
Scarani eq.\ 22 (QFT$_{n=2}$) & Prescribed matrix & max\_abs\_diff $=0$ (bit-identical) & \textbf{Yes} \\
Unitarity of QFT$_{n=2}$ & $F^\dagger F = \mathbf{1}$ & $\|F^\dagger F - \mathbf{1}\|_\infty = 0$ & \textbf{Yes} \\
\bottomrule
\end{tabular}

\section{Verdict}

\textbf{SPOT-CHECK} (per the QC brief's rule for survey-mode replications: no single headline number to reproduce, so we ran two representative algorithms + one direct exercise match, all of which hit their expected numbers).

Justification: Scarani's paper is a survey/pedagogical review, so there is no experimental table to reproduce directly.  Every quantitative object the paper does write down that we tested (Ex.\ 3.4 QFT matrix~eq.~22) reproduces bit-identically; every algorithm from the survey's citation frame that we exercised (Deutsch--Jozsa, referenced in Scarani's v2 addendum via Chuang~\emph{et~al.}\ \emph{Nature}~\textbf{393}) reproduces its textbook headline (\(P=1\) constant, \(P=0\) balanced); and Simon's algorithm --- the direct progenitor of Shor's period-finding (Scarani ref.~[2]) --- recovers all 5 tested hidden strings in $\le 4$ quantum rounds each.  Therefore the ``$R_u(\theta)$~+~CNOT'' universality thesis Scarani exposits actually delivers the promised quantum-advantage behaviour when driven by real Qiskit statevector simulation.

If Rick prefers a stricter vocabulary: this could also be labelled \textbf{PARTIAL}, because we did not exercise every one of Scarani's four exercises (only 3.4 was directly matched, and DJ/Simon go beyond the paper's own exercise set).

\section{Open Questions}
\begin{enumerate}
    \item[\textbf{Q1}] Scarani's Section~2 requires $\omega_1 - \omega_2 > 2\omega_c$ and $\ge 4\omega_c$ ``to work comfortably''; how does the CNOT infidelity in an actual noisy simulator (Aer with $T_1/T_2$ + amplitude-damping) scale with the ratio $(\omega_1-\omega_2)/\omega_c$?  In the ideal-gate statevector run we can never see this trade-off, so the paper's core NMR-physics claim (selectivity vs.\ speed) is untouched by our replication.  \emph{Basis:} the paper's selectivity/duration lower/upper bounds have no visible effect on our noise-free run.
    \item[\textbf{Q2}] The QFT statement in eq.~(21) is a \emph{DFT matrix}, not the standard efficient Coppersmith QFT circuit; when compiled to $R_u+\text{CNOT}$ on a real device (or on Qiskit's transpiled Aer), how many CNOTs are needed for $n=2,3,4$, and does the survey's implicit ``one-liner'' definition mislead a reader as to gate cost?  \emph{Basis:} we reproduced the $n{=}2$ matrix exactly, but the paper never counts gates.
    \item[\textbf{Q3}] Scarani presents his exercises with $\{|+\rangle,|-\rangle\}$ notation and a physicist's clockwise-positive rotation convention.  Multiple modern textbooks and Qiskit itself use the opposite sign convention for $R_x$.  Which of Scarani's exercise circuits produce a sign-flipped output under Qiskit's convention, and does re-solving Ex.\ 3.1 (GHZ) under Qiskit's convention require inverting any of the $R_y(\pi/4)$?  \emph{Basis:} our replication used the standard Qiskit basis $H,\text{CNOT}$ rather than Scarani's $R_x,R_y$ literally, so a full literal port of his exercise set would need this convention audit.
    \item[\textbf{Q4}] In Simon's algorithm, our implementation used the ``textbook'' oracle (copy $x$ into $y$, then XOR $s$ conditioned on the first set bit of $s$).  For real NMR-style spins (Scarani's frame), what is the shallowest oracle depth that still implements a valid 2-to-1 $f$ with hidden period $s$, and how does depth scale with $|s|_H$ (Hamming weight)?  \emph{Basis:} our max Simon depth of $\le 4$ rounds all recovered $s$, but oracle depth was never a variable we swept.
    \item[\textbf{Q5}] Scarani ends by conceding that quantum error correction is ``for specialists'' and skips it.  What fraction of the algorithmic circuits used in this replication (DJ~$n{=}3$, Simon~$n{=}3$) would fail on a physical device at typical 2026-era single/two-qubit gate infidelities ($p_1{\sim}10^{-3}$, $p_2{\sim}5\!\times\!10^{-3}$)?  Answering this would let a follow-up show a reader when the survey's neglect of QEC starts to matter as $n$ grows.  \emph{Basis:} our run was noise-free and finished in $<3$~s; the paper leaves a big instructive gap between ``it works in principle'' and ``it works on a real device.''
\end{enumerate}

\end{document}
